\section{Expectation and Variance of Adjudication-Induced Curvature}

Let $\Delta I = \ln 2$ denote a single adjudication event localized around $x$. The induced curvature adjustment is modelled as
\[
\Delta R_{\mu\nu}(x) = 8 \pi G \, k_B T \, \ln 2 \, \nabla_\mu \nabla_\nu \Phi(x; \sigma)
\]
where $\Phi$ solves the Green equation $\Box \Phi = \delta^{(4)}(x - x_0)$ with regulator $\sigma$ determined by the lattice spacing and background metric. The PT sources ensure that the event distribution is Poisson with rate density $\rho_{PT}(x)$.

## Expectation

Taking the expectation over independent events with density $\rho_{PT}$,
\[
\mathbb{E}[\Delta R_{\mu\nu}(x)] = 8 \pi G \, k_B T \, \ln 2 \int_X \rho_{PT}(x') \, \nabla_\mu \nabla_\nu \Phi(x - x') \, d^4 x'.
\]
Quarter-Lock invariance forces the reaction term to lie normal to the QL plane, yielding
\[
\mathbb{E}[\Delta R_{\mu\nu}] = 8 \pi G ( T^{(\Psi)}_{\mu\nu} + C_{\mu\nu} )
\]
matching the Einstein tensor recorded in TE1.C slow-roll diagnostics. A reference to `TE_1.C_RQG/results/phase1_summary.json` will supply the numerical targets.

## Variance

Because adjudication events follow a Poisson process, variance accumulates linearly with the rate:
\[
\mathrm{Var}[\Delta R_{\mu\nu}(x)] = ( 8 \pi G \, k_B T \, \ln 2 )^2 \int_X \rho_{PT}(x') \left( \nabla_\mu \nabla_\nu \Phi(x - x') \right)^2 d^4 x'.
\]
For homogeneous backgrounds the integral simplifies to $\rho_{PT} \cdot \kappa_{\mu\nu}$, where $\kappa_{\mu\nu}$ is a geometry-dependent constant computed from the lattice Green function. This yields
\[
\mathrm{Var}[\Delta R_{\mu\nu}] = ( 8 \pi G \, k_B T \, \ln 2 )^2 \, \rho_{PT} \, \kappa_{\mu\nu}.
\]
To calibrate $\kappa_{\mu\nu}$ we will measure the variance in the lattice simulator and compare against this prediction, using TE1.C temperature settings (`configs/spectra_slow_roll.yaml`).

## Next Steps
1. Specify the regulated Green function and compute $\kappa_{\mu\nu}$ analytically for the FRW background.
2. Cross-reference PT rate density $\rho_{PT}$ from TE1.R closure data (`results/pt_selector/*`).
3. Integrate this section into `proofs/RQG_quantization.tex` alongside corollaries.

\subsection{Numerical Validation}

Using the lattice configuration in `configs/lattice_simulation.yaml`, we
executed the Monte Carlo routine (`analysis/rqg_lattice.py`) with 10,000
adjudication events and variance calibration informed by
`results/kappa_baseline.json`. The recorded tensors (`results/lattice_events.npz`,
`results/lattice_summary.json`) satisfy
\begin{align*}
\langle\Delta R_{ij}\rangle &\approx 4.60\times10^{-32}, \\[-2pt]
\operatorname{Var}(\Delta R_{ij}) &\approx 2.39\times10^{-66},
\end{align*}
in agreement with the analytic prediction
$\operatorname{Var}(\Delta R_{ij})=(8\pi G k_B T \ln 2)^2 \rho_{\mathrm{PT}}\,\kappa_{ij}$.

Applying the spectrum fitting routine (`analysis/spectrum_fit.py` with
`step-init = 4.3744235775\times10^{-32}`) yields the discrete quantum increment
\[
\Delta R_{\text{fit}} =(4.374\pm0.022)\times10^{-32},
\]
confirming that curvature excitations occur in uniform steps set by the
adjudicative energy scale. The corresponding histogram and detailed statistics
are stored in `results/spectrum_summary.json`.
